\subsection*{两个量都趋零，比例仍可能完全不同}
$x/x$、$x^2/x$ 与 $x/x^2$ 的分子分母都趋零，比值却分别趋向一、零，或在零附近无界。零除以零不是一种答案，它只是提示直接代入没有揭示相对消失速度。需要进一步比较两个量的一阶或更高阶变化。

若 $f(a)=g(a)=0$，且两者可导、$g'(a)\ne0$，局部近似给出 $f(a+h)\approx f'(a)h$ 与 $g(a+h)\approx g'(a)h$，因此比例接近导数之比。这个简单版本提供了直觉，但一般洛必达法则可以处理更复杂的去心区间情况，需要明确的条件。

\subsection*{法则是有条件的工具}
下面列出的条件要求相关区间可微、分母导数不为零，并且导数比值具有极限。使用时先判断原式的未定形式，再核对这些条件。仅仅看到一个分式，不意味着可以把分子分母各求一次导数。

例如 $(x+1)/(x+2)$ 在零附近没有未定性，直接连续代入得 $1/2$；随意求导反而得到一。又如 $(x^2-1)/(x-1)$ 在一处，可以因式分解直接求极限，没必要把简单结构藏到一般定理里。

如果一次求导后仍是未定式，再用一次也要重新核对条件。洛必达法则的假设不满足，并不能反推原极限不存在，只能说明这个工具暂时不给结论。解题时应先看能否约分、换元或利用已知局部近似，再考虑更一般的方法。
